% =====================================================================
%  02-v2  (new-born paper #2)  --  SKELETON, started 2026-10-30
%  Source of all mathematics: ../02_fixed-h-mcfg/main.tex  ("v1", frozen).
%  Read INSTRUCTIONS_FOR_AI.md before editing.  Every \port{...} box says
%  what to move here from v1 (by label) and how to shorten it.
%  Nothing in this file is a new mathematical claim: Theorems A--C restate
%  results proved in v1 (labels given in the boxes and in the instructions).
% =====================================================================
\documentclass[11pt]{article}

\usepackage[T1]{fontenc}
\usepackage{lmodern}
\usepackage{microtype}
\usepackage{amsmath,amssymb,amsthm,mathtools}
\numberwithin{equation}{section}
\usepackage{enumitem}
\usepackage{array}
\usepackage{tabularx}
\usepackage{geometry}
\geometry{margin=1in}
\usepackage{cite}
\usepackage{hyperref}
\usepackage{url}
\usepackage{xcolor}
\hypersetup{hidelinks}
\setlength{\emergencystretch}{3em}

\theoremstyle{plain}
\newtheorem{theorem}{Theorem}[section]
\newtheorem{lemma}[theorem]{Lemma}
\newtheorem{proposition}[theorem]{Proposition}
\newtheorem{corollary}[theorem]{Corollary}
\newtheorem{mainthm}{Theorem}
\renewcommand{\themainthm}{\Alph{mainthm}}
\theoremstyle{definition}
\newtheorem{definition}[theorem]{Definition}
\theoremstyle{remark}
\newtheorem{remark}[theorem]{Remark}
\newtheorem{openproblem}[theorem]{Open Problem}

% --- macros kept identical to v1 so that ported text compiles unchanged ---
\newcommand{\Wit}{\mathsf{W}}
\newcommand{\Cmcf}[2]{\mathcal{C}^{\mathrm{mcf}}_{#1,#2}}
\newcommand{\D}{\mathcal{D}}
\newcommand{\out}{\mathrm{out}}
\newcommand{\tuple}[1]{\vec{#1}}
\newcommand{\blank}{\square}
\newcommand{\eps}{\varepsilon}
\newcommand{\fan}{\operatorname{fan}}
\newcommand{\YSub}[1]{\mathsf{YSub}(#1)}

% --- skeleton-only macro: remove every \port box before submission ---
\newcommand{\port}[1]{\leavevmode\par\medskip\noindent\fbox{\parbox{\dimexpr\linewidth-2\fboxsep-2\fboxrule}{\small\color{blue!50!black}\textbf{PORT / TODO.} #1}}\par\medskip}

\title{Local Groups and the Fan-Out Dependence of Characteristic-Data Length\\
in Typed Distributional Learning of Multiple Context-Free Languages}
\author{Takayuki Kuriyama\\
Independent Researcher, Tokyo, Japan\\
\texttt{growup.kuriyama@gmail.com}}
\date{}

\begin{document}
\maketitle

\begin{abstract}
We study positive-data learning of multiple context-free languages in which
nonempty factors and tuples are compared only when their types under a fixed
morphism $s:\Sigma^+\to\mathsf S$ into a finite semigroup agree, continuing
the context-free case~\cite{kuriyamaCFG}.  We measure the total length of
characteristic samples relative to a presentation $G$, in terms of its size
$\|G\|$ and its rule thickness $\bar\theta_G$.  First, we determine where
characteristic data are cheap.  If the positive image of $s$ is a local group,
data of length $O(\|G\|^2\bar\theta_G)$ exist at every fan-out, and under every
typing some targets with asymptotically smallest presentations need
$\Omega(\|G\|\bar\theta_G)$, the exact order for the trivial typing.  At fan-out
one, local groups are exactly the pseudovarieties with uniform polynomial
data.  For individual typings onto $U_1$ the cost depends on the letters of
type $1$: two such letters force exponential data, and a single one already
forces data quadratic in $\bar\theta_G$ for presentations of logarithmic size.
Second, the boundary moves with fan-out: at fan-out two the typing onto $U_1$
with a single letter of type $1$ admits polynomial data, and the exponential
fan-out-one families are not tuple-substitutable under any fixed typing.  The
main tools are simultaneous shortening of interleaved tuple components with
ordered group products, and two-hole power transfer.
\end{abstract}

% =====================================================================
\section{Introduction}
\label{sec:intro}

\port{Write the motivation paragraph (about 10 lines): fixed finite-semigroup
typing as a learning bias \cite{kuriyamaCFG,CosteTyping2004}; Yoshinaka's
multidimensional substitutability \cite{Yoshinaka2011} as the untyped baseline;
the question of this paper: \emph{for which typings is characteristic data
polynomial, and does the answer depend on fan-out?}  Reuse sentences from v1
\S1 (lines 98--350) where possible, but keep the three-theorem storyline
below.  Target length of the whole Introduction: at most 4 pages.}

\paragraph{Setting.}
A \emph{typing} is a morphism $s:\Sigma^+\to\mathsf S$ onto a finite semigroup.
The class $\Cmcf{f}{s}$ consists of the languages whose nonempty part is a
multiple context-free language of fan-out at most $f$ and which are
$(f,s)$-tuple-substitutable: tuples of equal types that share one admissible
context (internal separators nonempty) have the same distribution.  For a
reduced good presentation $G$ we use its size $\|G\|$ and its rule thickness
$\bar\theta_G$, the largest length of a shortest sentence using a given rule.
Lower bounds are stated for every set-driven operator with characteristic
samples throughout the class; upper bounds for the reconstruction operator of
\cite{kuriyamaCFG} lifted to tuples.

\paragraph{Main results.}
Fix the alphabet, the fan-out $f$, a rule-rank bound $r_{\max}\ge2$ and the
typing $s$.

\begin{mainthm}[Local groups and the linear factor]
\label{thm:A}
\begin{enumerate}[label=\textup{(\roman*)},leftmargin=*]
\item If the positive image of $s$ is a local group, every
$L\in\Cmcf{f}{s}$ has characteristic data of total length
$O(\|G\|^2\bar\theta_G)$ for every reduced good presentation $G$.
\item For every $f\ge1$ and every typing $s$ there are finite targets in
$\Cmcf{f}{s}$, with presentations of asymptotically smallest size, whose
characteristic samples all have length $\Omega(\|G\|\bar\theta_G)$.
\item Hence for the trivial typing the order $\Theta(\|G\|\bar\theta_G)$ is exact
at every fan-out.
\end{enumerate}
\end{mainthm}

\begin{mainthm}[The boundary at fan-out one]
\label{thm:B}
Let $f=1$.
\begin{enumerate}[label=\textup{(\roman*)},leftmargin=*]
\item A pseudovariety of finite semigroups admits uniform polynomial
characteristic-data length iff it is contained in $\mathbf{LG}$.
\item If some copy of $U_1$ in $\mathsf S$ is realized by two letters of its
identity type and one letter of its absorbing type, the class needs
exponential characteristic data.
\item For $\Sigma=\{\mathtt a,\mathtt b\}$ and $\mathtt a\mapsto1$,
$\mathtt b\mapsto\zeta$ onto $U_1$, there are targets with $\|G\|=O(\log\bar\theta_G)$
whose characteristic samples exceed $\bar\theta_G^{\,2}/4$; in particular no bound
$p(\|G\|)\,\bar\theta_G$ with $p$ a polynomial holds.
\end{enumerate}
\end{mainthm}

\begin{mainthm}[The boundary moves at fan-out two]
\label{thm:C}
\begin{enumerate}[label=\textup{(\roman*)},leftmargin=*]
\item For the typing of Theorem~\ref{thm:B}\textup{(iii)} and $f=2$, every target
has characteristic data of length $O(\|G\|^3\bar\theta_G^{\,5})$.
\item The exponential family behind Theorem~\ref{thm:B}\textup{(i)} is not
$(f,s)$-tuple-substitutable for any fixed finite typing and any $f\ge2$.
\end{enumerate}
\end{mainthm}

\port{\textbf{Where each part is proved in v1} (port statement + proof, keep
labels):
A(i) = \texttt{thm:local-group-polydata} (with \texttt{thm:window-group-polydata}
as the special case);
A(ii) = \texttt{prop:linear-factor-all-fanout} (the three-letter
\texttt{prop:linear-factor-sharp} becomes a remark);
A(iii) = trivial typing: $t_{\mathrm{typ}}\le\bar\theta_G$ and
\texttt{thm:data-from-typed-thickness};
B(i) = \texttt{thm:fanout-one-data} (+ \texttt{prop:bracketed-family});
B(ii) = \texttt{cor:typing-level-obstruction} (+ \texttt{prop:three-letter-family});
B(iii) = \texttt{prop:binary-u1-quadratic};
C(i) = \texttt{thm:binary-u1-polydata};
C(ii) = \texttt{thm:doubling-no-lift} (+ \texttt{lem:power-transfer}).
\textbf{If the open problem S1 below is resolved, Theorem~\ref{thm:B}(iii) and
Theorem~\ref{thm:C} change: see INSTRUCTIONS\_FOR\_AI.md \S5.}}

\paragraph{The remaining question.}
Theorems~\ref{thm:B}(iii) and~\ref{thm:C}(i) bracket one typing: at fan-out one
it needs data quadratic in $\bar\theta_G$, at fan-out two it needs only
polynomial data.  Whether its data at fan-out one are polynomial is the
central open problem (Open Problem~\ref{op:S1}).  A negative answer would show
that raising the fan-out can make learning strictly cheaper; a positive answer
would classify all typings onto $U_1$ at fan-out one by their number of
neutral letters.

\paragraph{Techniques.}
\port{Half a page.  (1) Owner-wise simultaneous replacement of interleaved
tuple components preserving ordered products in a finite group, with
Ramsey-type order homogenization for noncommutative groups
(\texttt{prop:block-replacement-interface}, App.~A of v1); this has no
fan-out-one analogue.  (2) Thinning and padding of owners for local groups
(\texttt{lem:lg-buffered-owners}).  (3) Distribution representatives instead
of generated witnesses (\texttt{cor:distribution-representatives}).
(4) Two-hole power transfer and Parikh saturation
(\texttt{lem:power-transfer}, \texttt{lem:app-u1-transfer}, \texttt{lem:app-u1-deep}).
(5) Marker families whose every subset stays in the class
(\texttt{prop:linear-factor-all-fanout}).}

\paragraph{Relation to prior work and to the companion paper.}
\port{Port v1 ``Relation to prior work'' and Table \texttt{tab:companion-overlap}
(v1 lines ~250--320), shortened.  Must state clearly: what is new relative to
\cite{kuriyamaCFG} (local groups instead of windows and groups; exactness at
fan-out one; typing-level results; all fan-out results), and that qualitative
learning is credited to Yoshinaka \cite{Yoshinaka2011}.}

\paragraph{Organization.}
\port{One paragraph, written last.}

% =====================================================================
\section{Preliminaries}
\label{sec:prelim}
\port{Port v1 \S2 verbatim (lines 352--647): positive-data learning,
\texttt{def:characteristic-sample}, typings, \texttt{def:prefix-suffix-typing},
\texttt{def:standard-mcfg}, \texttt{def:grammar-properties},
\texttt{def:tuple-occurrence}, \texttt{def:tuple-substitutable}.
Decide whether \texttt{prop:abelian-coset-filter} is still needed (grep its
uses); if not, park it.}

% =====================================================================
\section{Reconstruction and Typed Witnesses}
\label{sec:reconstruction}
\port{Port v1 \S3 (lines 648--962) but \textbf{shorten}: keep definitions
(\texttt{def:reconstruction}, \texttt{def:output-refinement},
\texttt{def:target-canonical}), the statements of
\texttt{prop:reconstruction-soundness}, \texttt{prop:completeness},
\texttt{prop:good-normalization}, \texttt{thm:exact-reconstruction},
\texttt{thm:poly}, \texttt{cor:gold-explanatory}; move proofs that only adapt
\cite{kuriyamaCFG,Yoshinaka2011} to the appendix.  Credit Yoshinaka.}

% =====================================================================
\section{Characteristic Data from Thickness and Representatives}
\label{sec:data}
\port{Port v1 \S4.1--4.2 (lines 963--1541): \texttt{lem:thickness-equivalence},
\texttt{lem:witness-length-template}, \texttt{lem:typed-thickness-equivalence},
\texttt{thm:data-from-typed-thickness}, \texttt{prop:rare-type},
\texttt{cor:typed-sandwich}, \texttt{cor:distribution-representatives}.
Reduce \texttt{prop:typed-thickness-exponential} to a remark (it limits the
witness method, not the data).  \textbf{Do not} put the linear-factor
propositions here: they go to \S\ref{sec:linear}.}

% =====================================================================
\section{Upper Bound: Local Groups at Every Fan-Out}
\label{sec:upper}
\port{Theorem~\ref{thm:A}(i).  Port v1 \S4.3--4.5 (lines 1542--1868):
\texttt{lem:heavy-node-count}, \texttt{prop:block-replacement-interface},
\texttt{lem:marked-heavy-window-group}, \texttt{thm:window-group-polydata},
\texttt{lem:mcfg-li-group-window}, \texttt{def:local-group-typing},
\texttt{thm:local-group-polydata}.  State at the start of the section that
the interleaving of components across subtrees is the MCFG-specific
difficulty.}

% =====================================================================
\section{Lower Bound: The Linear Factor at Every Fan-Out}
\label{sec:linear}
\port{Theorem~\ref{thm:A}(ii)--(iii).  Main statement:
\texttt{prop:linear-factor-all-fanout} (v1 lines 1423--1461, proof App.\ I of
v1).  \texttt{prop:linear-factor-sharp} (three letters, $f=1$) becomes a remark
with its short proof or is cited as the case that motivated the marker
construction.  Add the one-line derivation of A(iii).}

% =====================================================================
\section{The Boundary at Fan-Out One}
\label{sec:fanout-one}
\port{Theorem~\ref{thm:B}.  Port: \texttt{prop:bracketed-family},
\texttt{def:uniform-polydata}, \texttt{thm:fanout-one-data},
\texttt{rem:fanout-one-data} (rewrite (b) as the typing-level summary),
\texttt{prop:three-letter-family}, \texttt{cor:typing-level-obstruction}
(v1 lines 1869--2051), and \texttt{prop:binary-u1-quadratic} (v1 lines
2528--2556, proof App.\ J of v1).  End the section with the typing-level
trichotomy for typings onto $U_1$ (two neutral letters: exponential; one:
quadratic lower bound, polynomiality open; local groups: polynomial).}

% =====================================================================
\section{The Boundary Moves at Fan-Out Two}
\label{sec:fanout-two}
\port{Theorem~\ref{thm:C}.  Port \texttt{lem:power-transfer},
\texttt{thm:doubling-no-lift} (v1 lines 2387--2459),
\texttt{thm:binary-u1-polydata} with its outline and
\texttt{rem:binary-u1-scope} (v1 lines 2460--2527), proof App.\ H of v1.
Shorten \texttt{lem:app-u1-runs} to one sentence in the scope remark unless it
is needed elsewhere.}

% =====================================================================
\section{Conclusion and Open Problems}
\label{sec:conclusion}
\port{Rewrite around Theorems A--C.  Keep the ``Tightness of the exponents''
discussion (v1 lines ~2600--2640) but shortened.}

\begin{openproblem}[The single-neutral-letter typing at fan-out one]
\label{op:S1}
For $\Sigma=\{\mathtt a,\mathtt b\}$, $\mathtt a\mapsto1$, $\mathtt b\mapsto\zeta$ onto
$U_1$, fan-out one and $r_{\max}\ge2$: is the characteristic-data length
bounded by a polynomial in $\|G\|$ and $\bar\theta_G$?
\end{openproblem}

\begin{openproblem}[Pseudovarieties at higher fan-out]
\label{op:fanout-two}
\port{Port v1 \texttt{op:fanout-two}: for $f\ge2$, which pseudovarieties outside
$\mathbf{LG}$ admit uniform polynomial data?  Note that every pseudovariety not
contained in $\mathbf{LG}$ contains $U_1$, so a superpolynomial lower bound for
one typing onto $U_1$ at fan-out two would make $\mathbf{LG}$ the boundary at
fan-out two as well.}
\end{openproblem}

\section*{Declaration of generative AI and AI-assisted technologies in the manuscript preparation process}
During preparation of this work, the author used ChatGPT (OpenAI) and Claude
(Anthropic) to assist with organization of the manuscript, expression and
translation of text, cross-checks of references and notation, and editing of
\LaTeX{} source. After using these tools, the author reviewed and edited the
material as needed and takes full responsibility for the content of the publication.

% =====================================================================
\appendix
\section{Appendices to be ported}
\port{Port, in this order, from v1: A block replacement in finite groups
(\texttt{def:block-family} \ldots\ \texttt{lem:order-homogenize});
B local-group coordinates (\texttt{lem:local-group-morphism} \ldots\
\texttt{lem:local-group-shortest}, incl.\ \texttt{lem:lg-buffered-owners},
\texttt{thm:lg-characterization}, \texttt{prop:li-g-strict});
C qualitative reconstruction; D typed witnesses; E good normalization;
F bracketed doubling family; G three-letter family;
H binary typing at fan-out two (\texttt{app:binary-u1});
I linear factor at every fan-out (\texttt{lem:app-lf-path},
\texttt{lem:app-lf-descent}; \textbf{park} \texttt{prop:unary-cyclic-linear});
J quadratic data at fan-out one (\texttt{app:binary-u1-quadratic}).
\textbf{Park} (move to \texttt{parked/} with a note): v1 \S5 Rule Rank One
(\texttt{thm:rankone-polydata}), v1 \S6.1--6.2 comparison with
multidimensional substitutability and typing separations
(\texttt{def:yoshinaka-pd-substitutability}, \texttt{lem:sl2-boundary-control},
\texttt{prop:feature-cross-typed}, \texttt{lem:flag-separation},
\texttt{prop:lg-needed-target}, \texttt{prop:nonabelian-needed-target}),
\texttt{prop:unary-cyclic-linear}.}

\begin{thebibliography}{99}

\bibitem{Gold1967}
E.~M.~Gold.
\newblock Language identification in the limit.
\newblock {\em Information and Control}, 10(5):447--474, 1967.

\bibitem{Angluin1980}
D.~Angluin.
\newblock Inductive inference of formal languages from positive data.
\newblock {\em Information and Control}, 45(2):117--135, 1980.

\bibitem{ClarkEyraud2007}
A.~Clark and R.~Eyraud.
\newblock Polynomial identification in the limit of substitutable context-free
languages.
\newblock {\em Journal of Machine Learning Research}, 8:1725--1745, 2007.

\bibitem{deLaHiguera1997}
C.~de~la~Higuera.
\newblock Characteristic sets for polynomial grammatical inference.
\newblock {\em Machine Learning}, 27:125--138, 1997.

\bibitem{CosteTyping2004}
F.~Coste, D.~Fredouille, C.~Kermorvant, and C.~de~la~Higuera.
\newblock Introducing domain and typing bias in automata inference.
\newblock In \emph{Grammatical Inference: Algorithms and Applications -- ICGI 2004},
LNCS 3264, pp.~115--126, Springer, 2004.

\bibitem{Takada1988}
Y.~Takada.
\newblock Grammatical inference for even linear languages based on control sets.
\newblock \emph{Information Processing Letters}, 28(4):193--199, 1988.
\newblock doi:10.1016/0020-0190(88)90208-6.

\bibitem{Takada1995}
Y.~Takada.
\newblock Learning formal languages based on control sets.
\newblock In K.~P. Jantke and S.~Lange (eds.), \emph{Algorithmic Learning
for Knowledge-Based Systems}, LNCS 961, pp.~316--339, Springer, 1995.

\bibitem{Clark2013}
A.~Clark.
\newblock Learning trees from strings: A strong learning algorithm for some context-free grammars.
\newblock \emph{Journal of Machine Learning Research}, 14:3537--3559, 2013.

\bibitem{CosteGaretNicolas2012}
F.~Coste, G.~Garet, and J.~Nicolas.
\newblock Local substitutability for sequence generalization.
\newblock In {\em Proceedings of the 11th International Conference on
Grammatical Inference}, JMLR Workshop and Conference Proceedings 21,
pp.~97--111, 2012.

\bibitem{EyraudHeinzYoshinaka2016}
R.~Eyraud, J.~Heinz, and R.~Yoshinaka.
\newblock Efficiency in the identification in the limit learning paradigm.
\newblock In J.~Heinz and J.~M.~Sempere (eds.), {\em Topics in Grammatical
Inference}, pp.~25--46. Springer, 2016.
\newblock doi:10.1007/978-3-662-48395-4\_2.

\bibitem{CosteNicolas2019}
F.~Coste and J.~Nicolas.
\newblock Learning local substitutable context-free languages from positive
examples in polynomial time and data by reduction.
\newblock In {\em Proceedings of the 14th International Conference on
Grammatical Inference}, Proceedings of Machine Learning Research 93,
pp.~155--168, 2019.

\bibitem{Yoshinaka2008}
R.~Yoshinaka.
\newblock Identification in the limit of $k,l$-substitutable context-free
languages.
\newblock In {\em Grammatical Inference: Algorithms and Applications},
LNCS 5278, pp.~266--279. Springer, 2008.

\bibitem{Yoshinaka2009}
R.~Yoshinaka.
\newblock Learning mildly context-sensitive languages with multidimensional
 substitutability from positive data.
\newblock In {\em Algorithmic Learning Theory}, LNCS 5809, pp.~278--292.
 Springer, 2009.
\newblock doi:10.1007/978-3-642-04414-4\_24.

\bibitem{Yoshinaka2011}
R.~Yoshinaka.
\newblock Efficient learning of multiple context-free languages with
multidimensional substitutability from positive data.
\newblock {\em Theoretical Computer Science}, 412(19):1821--1831, 2011.
\newblock doi:10.1016/j.tcs.2010.12.058.

\bibitem{Yoshinaka2015General}
R.~Yoshinaka.
\newblock General perspective on distributionally learnable classes.
\newblock In \emph{Proceedings of the 14th Meeting on the Mathematics
of Language (MoL 2015)}, pp.~87--98.
Association for Computational Linguistics, 2015.
\newblock doi:10.3115/v1/W15-2308.

\bibitem{Yoshinaka2010}
R.~Yoshinaka.
\newblock Polynomial-time identification of multiple context-free languages
from positive data and membership queries.
\newblock In {\em Grammatical Inference: Theoretical Results and Applications},
LNCS 6339, pp.~230--244. Springer, 2010.
\newblock doi:10.1007/978-3-642-15488-1\_19.

\bibitem{ClarkYoshinaka2016}
A.~Clark and R.~Yoshinaka.
\newblock Distributional learning of context-free and multiple context-free grammars.
\newblock In J.~Heinz and J.~M.~Sempere (eds.), \emph{Topics in Grammatical
Inference}, pp.~143--172. Springer, 2016.
\newblock doi:10.1007/978-3-662-48395-4\_6.

\bibitem{ClarkYoshinaka2014}
A.~Clark and R.~Yoshinaka.
\newblock Distributional learning of parallel multiple context-free grammars.
\newblock {\em Machine Learning}, 96(1--2):5--31, 2014.
\newblock doi:10.1007/s10994-013-5403-2.

\bibitem{OatesEtAl2006}
T.~Oates, T.~Armstrong, L.~Becerra-Bonache, and M.~Atamas.
\newblock Inferring grammars for mildly context-sensitive languages in polynomial-time.
\newblock In {\em Grammatical Inference: Algorithms and Applications -- ICGI 2006},
LNCS 4201, pp.~137--147. Springer, 2006.

\bibitem{KanazawaYoshinaka2023}
M.~Kanazawa and R.~Yoshinaka.
\newblock Extending distributional learning from positive data and membership queries.
\newblock In {\em Proceedings of the 16th International Conference on Grammatical Inference},
Proceedings of Machine Learning Research 217, pp.~8--22, 2023.

\bibitem{DuncanEtAl2026Substitution}
A.~Duncan, M.~Elder, L.~Frenkel, and M.~Lyu.
\newblock A substitution lemma for multiple context-free languages.
\newblock \emph{International Journal of Algebra and Computation},
2026 (online publication).
\newblock doi:10.1142/S0218196726500463.

\bibitem{RambowSatta1999}
O.~Rambow and G.~Satta.
\newblock Independent parallelism in finite copying parallel rewriting systems.
\newblock {\em Theoretical Computer Science}, 223(1--2):87--120, 1999.
\newblock doi:10.1016/S0304-3975(97)00190-4.

\bibitem{kuriyamaCFG}
T.~Kuriyama.
\newblock Distributional learning of context-free languages under fixed
finite-semigroup typing.
\newblock \url{https://arxiv.org/abs/1409.6247v5}.

\bibitem{Pin1986}
J.-E. Pin.
\newblock {\em Varieties of Formal Languages}.
\newblock North Oxford Academic, London, and Plenum, New York, 1986.

\bibitem{Costa2002LI}
J.~C.~Costa.
\newblock Some pseudovariety joins involving locally trivial semigroups.
\newblock \emph{Semigroup Forum}, 64(1):12--28, 2002.
\newblock doi:10.1007/s002330010060.

\bibitem{Costa2002LIG}
J.~C.~Costa.
\newblock Some pseudovariety joins involving groups and locally trivial semigroups.
\newblock In \emph{Semigroups, Algorithms, Automata and Languages},
pp.~341--348. World Scientific, 2002.

\bibitem{SekiEtAl1991}
H.~Seki, T.~Matsumura, M.~Fujii, and T.~Kasami.
\newblock On multiple context-free grammars.
\newblock {\em Theoretical Computer Science}, 88(2):191--229, 1991.
\newblock doi:10.1016/0304-3975(91)90374-B.

\bibitem{Kallmeyer2010}
L.~Kallmeyer.
\newblock {\em Parsing Beyond Context-Free Grammars}.
\newblock Springer, Berlin, Heidelberg, 2010.
\newblock doi:10.1007/978-3-642-14846-0.

\bibitem{Kanazawa2019Ogden}
M.~Kanazawa.
\newblock Ogden's lemma, multiple context-free grammars, and the control language hierarchy.
\newblock {\em Information and Computation}, 269:104449, 2019.
\newblock doi:10.1016/j.ic.2019.104449.

\bibitem{GomezRodriguezEtAl2009}
C.~G\'omez-Rodr\'iguez, M.~Kuhlmann, G.~Satta, and D.~Weir.
\newblock Optimal reduction of rule length in linear context-free rewriting systems.
\newblock In {\em Proceedings of Human Language Technologies: The 2009 Annual
Conference of the North American Chapter of the Association for Computational
Linguistics}, pp.~539--547. Association for Computational Linguistics, 2009.

\bibitem{GeroldingerHalterKoch2006}
A.~Geroldinger and F.~Halter-Koch.
\newblock {\em Non-Unique Factorizations: Algebraic, Combinatorial and
Analytic Theory}, volume 278 of {\em Pure and Applied Mathematics}.
\newblock Chapman \& Hall/CRC, Boca Raton, 2006.

\bibitem{AlmeidaEtAl2009}
J.~Almeida, S.~Margolis, B.~Steinberg, and M.~Volkov.
\newblock Representation theory of finite semigroups, semigroup radicals and
formal language theory.
\newblock {\em Transactions of the American Mathematical Society},
361(3):1429--1461, 2009.
\newblock doi:10.1090/S0002-9947-08-04712-0.

\bibitem{Straubing1985}
H.~Straubing.
\newblock Finite semigroup varieties of the form $\mathbf V*\mathbf D$.
\newblock {\em Journal of Pure and Applied Algebra}, 36:53--94, 1985.
\newblock doi:10.1016/0022-4049(85)90062-3.

\bibitem{Tilson1987}
B.~Tilson.
\newblock Categories as algebra: an essential ingredient in the theory of monoids.
\newblock {\em Journal of Pure and Applied Algebra}, 48:83--198, 1987.
\newblock doi:10.1016/0022-4049(87)90108-3.

\end{thebibliography}


\end{document}
